\documentclass[12pt,a4paper]{article}
\usepackage{amsmath,amssymb,amsthm}
\usepackage{geometry}
\usepackage{hyperref}
\usepackage{microtype}
\usepackage{booktabs}
\usepackage{array}
\usepackage{graphicx}
\geometry{margin=2.5cm}

\newtheorem{theorem}{Theorem}[section]
\newtheorem{lemma}[theorem]{Lemma}
\newtheorem{corollary}[theorem]{Corollary}
\newtheorem{definition}[theorem]{Definition}
\newtheorem{remark}[theorem]{Remark}
\newtheorem{observation}[theorem]{Observation}

\title{The Cascade Series --- Part IVa\\[0.5em]
\large The Standard Model from the Cascade:\\
Gauge Group, Symmetry Breaking, and Three Generations\\
from Bott Periodicity and Hairy Ball Zeros}
\author{RTAC}
\date{March 2026}

\begin{document}
\setlength{\emergencystretch}{3em}
\maketitle

\begin{abstract}
The cascade series tests one hypothesis: the infinite-dimensional unit ball, descended to
four dimensions, is indistinguishable from our universe. Papers~I--III~\cite{paperI,paperII,paperIII}
established the cosmological constant, quantum mechanics, and general relativity as
consequences of the cascade's geometry. This paper derives the Standard Model gauge group
$\mathrm{SU}(3)\times\mathrm{SU}(2)\times\mathrm{U}(1)$, its symmetry breaking pattern,
and the number of fermion generations from the same geometry, with no additional input.

The argument has three components. (1)~\emph{Gauge group from Bott periodicity and Adams'
theorem.} The Clifford algebra $\mathrm{Cl}(1,d-1)$ has complex minimal spinors when
$d\bmod 8\in\{4,5,6\}$. The cascade tower contains a second complex window at $\{12,13,14\}$,
the Bott mirror of the spacetime window $\{4,5,6\}$. Adams' theorem determines the gauge
group dimensions via the Radon--Hurwitz number: $\rho(12)-1=3$ independent nowhere-zero
vector fields on $S^{11}$ gives $N_c=3$ colours; $\rho(14)-1=1$ on $S^{13}$ gives
$\dim\mathrm{U}(1)=1$. The 12 layers between $d_0=7$ and $d_1=19$ match the 12 generators
of $\mathrm{SU}(3)\times\mathrm{SU}(2)\times\mathrm{U}(1)$. Furthermore, $d=12$ is the
\emph{unique} dimension in $[5,d_1=19]$ where $\rho(d)-1=3$, so the gauge window is
forced, not chosen. (2)~\emph{Symmetry breaking from the hairy ball theorem.} The hairy
ball theorem forces every tangent vector field on even-dimensional spheres to have a zero.
A stronger group-theoretic statement holds: no connected compact Lie group can act freely
on an even-dimensional sphere (Lefschetz fixed-point theorem). In the gauge window:
$d=12$ operates on $S^{11}$ (odd, no forced zero $\to$ $\mathrm{SU}(3)$ unbroken), $d=13$
on $S^{12}$ (even, Lefschetz obstruction $\to$ $\mathrm{SU}(2)$ broken), $d=14$ on
$S^{13}$ (odd, no forced zero $\to$ $\mathrm{U}(1)$ unbroken). The Standard Model
symmetry breaking pattern is a topological theorem. The Higgs mass is the geodesic
distance from the zero to the VEV on $S^{12}$: $m_H/m_W=\pi/2$, giving $m_H=126.3$ GeV
(observed 125.25 GeV, deviation 0.80\%). (3)~\emph{Three generations from the phase
transition at $d_1=19$.} A fermion generation requires a complex Dirac layer ($d\bmod 8=5$)
with a hairy ball zero ($d$ odd). These occur at $d=5,13,21,29,\ldots$ The first threshold
$d_1=19$ is a phase transition in the decay rate $p(d)$: subcritical for $d<d_1$,
supercritical for $d>d_1$. Three generation layers ($d=5,13,21$) straddle this transition;
the fourth Bott fermion layer at $d=29$ is deep in the supercritical regime and suppressed
from the charged-fermion spectrum by $\sim 289$ (Theorem~\ref{thm:generations}), with its
residual role as the heaviest-neutrino source derived in Part~IVb.

The companion paper~\cite{paperIVb} derives the fermion mass spectrum, gauge coupling
constants, and precision predictions from these structural results.
\end{abstract}

\tableofcontents
\newpage

\section{Division of Labour}

Papers~I--III of the cascade series derive, from the cascade's geometry alone: a geometric
invariant $I=1.0990\times 10^{-120}$ matching the cosmological constant~\cite{paperI}; the
structural framework of quantum mechanics~\cite{paperII}; and general relativity with
$\Lambda=I$, spacetime dimension $d=4$, and Lorentzian signature~\cite{paperIII}. This
paper and its companion~\cite{paperIVb} complete the programme by deriving the Standard
Model.

The cascade provides geometric content; classical mathematical theorems provide the
uniqueness. The new ingredients in this paper are Bott periodicity (which organises the
Clifford algebra classification by dimension), the hairy ball theorem and Lefschetz
fixed-point theorem (which determine which gauge symmetries are broken), Adams' theorem
with its uniqueness property (which fixes the gauge window and colour count), and the
first threshold $d_1=19$ as a phase transition (which limits the number of fermion
generations).

\begin{center}
\resizebox{\textwidth}{!}{%
\begin{tabular}{lll}
\toprule
Paper & Cascade provides & Uniqueness provides \\
\midrule
I~\cite{paperI}    & $\Omega_d$ cascade, $I\approx 10^{-120}$ & --- (pure geometry) \\
II~\cite{paperII}  & Complex state space, propagator $e^{-iHt}$ & Sphere geometry (Born rule) \\
III~\cite{paperIII}& Metric, $\Lambda=I$, $d=4$, $(-,+,+,+)$ & Lovelock, Clifford, propagator \\
IVa (this)         & Gauge window, breaking pattern & Bott, Adams, Lefschetz \\
IVb~\cite{paperIVb}& Mass spectrum, couplings & Phase transition, cascade potential \\
\bottomrule
\end{tabular}%
}
\end{center}

\section{The Bott Mirror: Gauge Structure from Clifford Periodicity}

\subsection{Complex spinor windows in the cascade}

The Clifford algebra $\mathrm{Cl}(1,d-1)$ determines the spinor type at each spacetime
dimension $d$. The type of the minimal spinor of $\mathrm{Spin}(1,d-1)$ is determined by
$\mathrm{Cl}(1,d-1)$, following Bott periodicity with period 8:

\begin{center}
\begin{tabular}{ccccc}
\toprule
$d$ & $\mathrm{Cl}(1,d-1)$ structure & Min.\ spinor & Type & Complex? \\
\midrule
2 & $M_2(\mathbb{R})$ & 2 real & Majorana & No \\
3 & $M_2(\mathbb{R})\oplus M_2(\mathbb{R})$ & 2 real & Majorana & No \\
4 & $M_2(\mathbb{C})\otimes_{\mathbb{R}}M_2(\mathbb{R})$ & 2 complex & Weyl & Yes \\
5 & $M_4(\mathbb{C})$ & 4 complex & Dirac & Yes \\
6 & $M_4(\mathbb{C})\oplus M_4(\mathbb{C})$ & 4 complex & Weyl & Yes \\
7 & $M_8(\mathbb{R})$ & 8 real & Majorana & No \\
8 & $M_{16}(\mathbb{R})$ & 8 real & Majorana & No \\
9 & $M_{16}(\mathbb{R})\oplus M_{16}(\mathbb{R})$ & 16 real & Majorana & No \\
\bottomrule
\end{tabular}
\end{center}

Data from Lounesto~\cite{lounesto}, Lorentzian signature $(1,d-1)$. The minimal spinor
representation is irreducibly complex when $d\bmod 8\in\{4,5,6\}$; all other residues give
real (Majorana) spinors. The complex windows in the cascade tower are:

\begin{center}
\begin{tabular}{llll}
\toprule
Window & Dimensions & Spinor types & Role \\
\midrule
1 & $\{4,5,6\}$   & Weyl, Dirac, Weyl & Spacetime \\
2 & $\{12,13,14\}$ & Weyl, Dirac, Weyl & Gauge \\
3 & $\{20,21,22\}$ & Weyl, Dirac, Weyl & (see~\cite{paperIVb}) \\
\bottomrule
\end{tabular}
\end{center}

The second window $\{12,13,14\}$ is the Bott mirror of the spacetime window $\{4,5,6\}$.
It reproduces the same Weyl--Dirac--Weyl pattern exactly, shifted by one Bott period.
Between the two complex windows lie five real (Majorana) layers $d=7,8,9,10,11$, which
carry no complex spinor structure and support no quantum matter content visible to the $d=4$
observer.

\subsection{The self-dual crossing at \texorpdfstring{$d=12$}{d=12}}

The cascade's lapse function $N(d)=\sqrt{\pi}\cdot\Gamma((d+1)/2)/\Gamma((d+2)/2)$ is the
compactification radius at dimension $d$.

\begin{theorem}[Self-dual crossing --- near-miss observation]\label{thm:self-dual}
The cascade lapse $N(d)$ crosses the Kaluza--Klein self-dual reference
value $1/\sqrt{2}$ between integers $d=12$ and $d=13$:
\[
N(12) \;=\; \frac{10395\,\pi}{46080} \;=\; 0.70870
\;>\; \tfrac{1}{\sqrt{2}} \;=\; 0.70711
\;>\; 0.68198 \;=\; N(13).
\]
The integer $d=12$ sits $0.225\%$ above the self-dual value; the integer
$d=13$ sits $3.6\%$ below. The continuous solution of
$N(d^*) = 1/\sqrt{2}$ lies in the open interval $(12,13)$ and does not
coincide with either integer.
\end{theorem}

The value $1/\sqrt{2}$ enters this theorem as the Kaluza--Klein
\emph{reference} radius where winding and momentum modes are degenerate
in KK compactification on a circle; in the literature, a U(1) gauge
field enhances to SU(2) at that radius. The cascade refuses KK
reduction (Paper~I~\S3.2; Paper~II\,=\,III~\S5--6; Paper~III~\S12), and
the SU(3)/SU(2)/U(1) assignment to $d\in\{12,13,14\}$ is derived in
Theorem~\ref{thm:adams-unique} below from Adams' theorem and Bott
periodicity, not from the self-dual crossing. The crossing serves only
to name the reference value $1/\sqrt{2}$ and to locate the second
complex spinor window (which opens at $d=12$) in numerical proximity to
it; see Remark~\ref{rem:sp27-status} for Check-7 compliance.

\begin{center}
\begin{tabular}{cccc}
\toprule
$d$ & $N(d)$ & $1/\sqrt{2}$ & Status \\
\midrule
4   & 1.17810 & 0.70711 & Observer dimension \\
7   & 0.91429 & 0.70711 & Volume maximum $d_0$ \\
12  & 0.70870 & 0.70711 & $0.225\%$ above self-dual reference \\
13  & 0.68198 & 0.70711 & $3.6\%$ below self-dual reference \\
19  & 0.56755 & 0.70711 & First threshold $d_1$ \\
\bottomrule
\end{tabular}
\end{center}

The crossing of $N(d) = 1/\sqrt{2}$ occurs between $d=12$ and $d=13$,
with $d=12$ the closest integer above ($0.225\%$ deviation). The
integer $d=12$ is the first layer of the second complex spinor window
(by Bott periodicity), but its proximity to the self-dual value is a
numerical coincidence of the Gamma function's values at integer
arguments, not a structural identity: no cascade theorem forces
$N(12) = 1/\sqrt{2}$, and no cascade derivation of the Standard Model
gauge group depends on $d=12$ sitting exactly on the self-dual radius.

\begin{remark}[Derivation status of the self-dual crossing]\label{rem:sp27-status}
The self-dual crossing has three different derivation statuses
depending on what claim is being made:
\begin{enumerate}
\item \emph{Derived (Gamma-function computation).} The numerical value
$N(12) = 10395\pi/46080 = 0.70870$ is an exact Gamma-function
evaluation; similarly $N(13) = 0.68198$. These are not approximations.
\item \emph{Observed (not derived as equality).} The proximity
$|N(12) - 1/\sqrt{2}|/(1/\sqrt{2}) = 0.225\%$ is a numerical
near-miss, not a structural identity. No cascade theorem forces
$N(12) = 1/\sqrt{2}$; the Gamma function happens to take a value
close to $1/\sqrt{2}$ at the integer $d=12$.
\item \emph{Not imported (Check-7 compliance).} The Kaluza--Klein
gauge-enhancement mechanism at the self-dual radius (U(1)
$\to$ SU(2) when winding and momentum modes are degenerate) is
\emph{not} used as a cascade derivation step. The cascade's
SU(3)/SU(2)/U(1) assignment at $d\in\{12,13,14\}$ is derived in
Theorem~\ref{thm:adams-unique} from Adams' theorem (the unique
dimension in $[5,d_1]$ with $\rho(d)-1 = 3$) and Bott periodicity
(the Weyl--Dirac--Weyl pattern). The self-dual radius appears only
as a literature reference value naming $1/\sqrt{2}$; the cascade
never performs a KK reduction, a semiclassical compactification, or
a winding-momentum mode analysis.
\end{enumerate}
An earlier phrasing of this section --- ``the crossing therefore falls
exactly at the first layer of the second complex spinor window'' ---
overstated item~(2): the continuous solution of $N(d^*) = 1/\sqrt{2}$
lies strictly between $d=12$ and $d=13$, and $d=12$ sits $0.225\%$
above rather than \emph{at} the self-dual value. An earlier phrasing
---``this gauge enhancement mechanism arises in the cascade from the
Gamma function'' --- could be read as invoking the KK mechanism as an
intermediate step; the corrected framing above separates the KK
reference value (a \emph{name} for $1/\sqrt{2}$) from the KK mechanism
(which the cascade does not use).
\end{remark}

\subsection{Three gauge factors from three complex layers}

The gauge window $\{12,13,14\}$ contains three layers with complex spinor structure. The
Bott mirror assigns a spinor type and propagator phase to each:

\begin{center}
\begin{tabular}{ccccc}
\toprule
$d$ & Spinor & Phase & Gauge factor & Generators \\
\midrule
12 & Complex Weyl  & $e^{i\cdot 4\pi}=+1$ & $\mathrm{SU}(3)$ & 8 \\
13 & Complex Dirac & $e^{i\cdot 9\pi/2}=i$ & $\mathrm{SU}(2)$ & 3 \\
14 & Complex Weyl  & $e^{i\cdot 5\pi}=-1$  & $\mathrm{U}(1)$  & 1 \\
\bottomrule
\end{tabular}
\end{center}

The propagator phases are the fourth roots of unity: $\{+1,i,-1\}$. The phase at layer $d$
is $e^{i(d-4)\pi/2}$, accumulating $\pi/2$ per cascade step from the forced precession
(Theorem~6.1 of~\cite{paperII}). The Weyl--Dirac--Weyl pattern assigns chiral
representations to $d=12$ and $d=14$ and a non-chiral (Dirac) representation to
$d=13$. This is a Clifford-algebra classification of the spinor representation space at
each layer; it is distinct from the chirality structure of each gauge boson's coupling
to fermions in the Standard Model (which is vectorial for $\mathrm{SU}(3)_c$,
parity-violating for $\mathrm{SU}(2)_L$, and determined separately by the symmetry
breaking of Section~3 and the fermion generation structure of Section~4). The
assignment of $\mathrm{SU}(3)$, $\mathrm{SU}(2)$, $\mathrm{U}(1)$ to layers $d=12,13,14$
specifically is forced by Adams' theorem (Theorem~\ref{thm:adams-unique}), not by
matching Clifford type to coupling chirality.

The real phases ($+1$ at $d=12$, $-1$ at $d=14$) give parity-invariant propagator
coefficients at the $\mathrm{SU}(3)$ and $\mathrm{U}(1)_{\rm em}$ layers, consistent with
the observed vectorial character of both gauge interactions at the observer's frame.
The imaginary phase ($i$ at $d=13$) marks the $\mathrm{SU}(2)$ layer as the unique
propagator-odd layer in the gauge window; this is the layer at which the hairy ball
zero forces symmetry breaking (Section~3). The strong-sector CP phase $\theta_{\rm QCD}$
is independently set to zero by the cascade's topological classification (Part~IVb
Section~6), not by a leakage argument from the $d=13$ phase.

\subsection{Generator counting}

\begin{remark}[Generator count coincidence]\label{thm:generators}
Two independent computations yield the same integer:
\begin{itemize}
\item[(a)] \emph{Threshold spacing.} The integer separation between the
two threshold-bracketing equilibria of Part~0 is
$d_1 - d_0 = 19 - 7 = 12$ (Part~0~\cite{paperI}, Theorem~6.7).
\item[(b)] \emph{Gauge-algebra dimension.} The total dimension of the
Standard Model gauge algebra is
$\dim(\mathrm{SU}(3)) + \dim(\mathrm{SU}(2)) + \dim(\mathrm{U}(1)) =
8 + 3 + 1 = 12$, assembled from Adams' theorem applied to $d=12$ (the
unique layer in $[5,d_1]$ with $\rho(d)-1=3$, Theorem~\ref{thm:adams-unique}
below), the Clifford--Lefschetz analysis at $d=13$ (Section~\ref{sec:breaking}),
and Adams' theorem at $d=14$ ($\rho(14)-1=1$, Theorem~\ref{thm:adams}).
\end{itemize}
The equality $12 = 12$ is a \emph{numerical consistency check} between
two independently-derived counts --- Part~0's threshold spacing and the
cascade's gauge-window assembly --- not a structural identity forced by
either derivation alone. That the two counts coincide is a property of
the specific $(d_V,d_0,d_1,d_2)$ tower selected by the Gamma function; a
cascade with different thresholds would generically give different
integers for (a) and (b).
\end{remark}

The total rank of $\mathrm{SU}(3)\times\mathrm{SU}(2)\times\mathrm{U}(1)$
is $2+1+1=4$, which numerically equals the observer's spacetime dimension
$d=4$. The cascade does not independently derive that the gauge window's
total Lie-algebra rank \emph{must} equal the observer's spacetime
dimension; both values are fixed by their separate structural derivations
(Bott~$+$~Adams at the gauge window; Lovelock~$\cap$~Clifford for $d=4$
in Paper~III~\cite{paperIII}, Section~9), and the equality is a second
numerical coincidence of the same kind as the generator count. We note it
as a consistency check without claiming it is structurally forced.

\subsection{\texorpdfstring{$N_c=3$}{N\_c=3} from Adams' theorem}

\begin{theorem}[Adams' theorem applied to the gauge window]\label{thm:adams}
The maximum number of linearly independent nowhere-zero tangent vector fields on $S^{n-1}$
is $\rho(n)-1$, where the Radon--Hurwitz number is computed as follows: write $n=2^a\cdot m$
with $m$ odd; write $a=4q+r$ with $0\leq r\leq 3$; then $\rho(n)=8q+2^r$.
Applied to each layer of the gauge window:
\begin{center}
\begin{tabular}{ccccc}
\toprule
$d$ & Sphere & $\rho(d)$ & Max fields & Physical meaning \\
\midrule
12 & $S^{11}$ & 4 & 3 & $N_c=3$ colours \\
13 & $S^{12}$ & 1 & 0 & No nonvanishing field (broken) \\
14 & $S^{13}$ & 2 & 1 & $\dim\mathrm{U}(1)=1$ \\
\bottomrule
\end{tabular}
\end{center}
\end{theorem}
\begin{proof}
Write $n=2^a\cdot m$ with $m$ odd; write $a=4q+r$ with $0\leq r\leq 3$; then
$\rho(n)=8q+2^r$. For $n=12=2^2\cdot 3$: $a=2$, $q=0$, $r=2$, giving $\rho(12)=4$ and
$\rho-1=3$ independent nowhere-zero vector fields on $S^{11}$. For $n=13=2^0\cdot 13$:
$a=0$, $q=0$, $r=0$, giving $\rho(13)=1$ and $\rho-1=0$ fields. For $n=14=2^1\cdot 7$:
$a=1$, $q=0$, $r=1$, giving $\rho(14)=2$ and $\rho-1=1$ field.
\end{proof}

The Radon--Hurwitz number also confirms the spacetime window: at $d=4$, $S^3$ has
$\rho(4)-1=3$ independent vector fields, matching the 3 spatial dimensions. The same
topological invariant governs both the spacetime structure and the gauge structure, applied
at the two Bott mirrors.

\begin{theorem}[Uniqueness of the gauge window]\label{thm:adams-unique}
Among all dimensions $5\leq d\leq d_1=19$, the dimension $d=12$ is the \emph{unique}
dimension for which $\rho(d)-1=3$.
\end{theorem}
\begin{proof}
We compute $\rho(d)$ for each $d\in\{5,\ldots,19\}$ using the Radon--Hurwitz formula.
Write $d=2^a\cdot m$ with $m$ odd; write $a=4q+r$ with $0\leq r\leq 3$; then
$\rho(d)=8q+2^r$.

\begin{center}
\renewcommand{\arraystretch}{1.15}
\begin{tabular}{ccccc}
\toprule
$d$ & $a$ & $(q,r)$ & $\rho(d)$ & $\rho(d)-1$ \\
\midrule
5  & 0 & $(0,0)$ & 1 & 0 \\
6  & 1 & $(0,1)$ & 2 & 1 \\
7  & 0 & $(0,0)$ & 1 & 0 \\
8  & 3 & $(0,3)$ & 8 & 7 \\
9  & 0 & $(0,0)$ & 1 & 0 \\
10 & 1 & $(0,1)$ & 2 & 1 \\
11 & 0 & $(0,0)$ & 1 & 0 \\
\textbf{12} & 2 & $(0,2)$ & 4 & $\mathbf{3}$ $\checkmark$ \\
13 & 0 & $(0,0)$ & 1 & 0 \\
14 & 1 & $(0,1)$ & 2 & 1 \\
15 & 0 & $(0,0)$ & 1 & 0 \\
16 & 4 & $(1,0)$ & 9 & 8 \\
17 & 0 & $(0,0)$ & 1 & 0 \\
18 & 1 & $(0,1)$ & 2 & 1 \\
19 & 0 & $(0,0)$ & 1 & 0 \\
\bottomrule
\end{tabular}
\end{center}

Among all $d\in[5,19]$, exactly $d=12$ gives $\rho(d)-1=3$. Note that $d=8$ gives
$\rho(8)-1=7$, corresponding to the 7 imaginary units of the octonions; however,
$S^7$ is not a Lie group (Adams~\cite{adams60}: the only spheres that are $H$-spaces are
$S^0,S^1,S^3,S^7$, and $S^7$ is non-associative). The cascade correctly places $d=8$
as the Bott period layer, not a gauge layer. $d=12$ is unique.
\end{proof}

\begin{remark}
The gauge window is therefore not a selection but a uniqueness statement: the number
$N_c=3$ is a topological invariant of the cascade, not a free parameter.
\end{remark}

\begin{remark}[Gauge group uniqueness]
The gauge group assignment is derived, not assumed. The Lefschetz theorem
(Theorem~\ref{thm:lefschetz} below) determines which symmetries are broken.
Adams' theorem determines the dimensions of the unbroken gauge groups: $N_c=3$ for
$\mathrm{SU}(3)$ and 1 for $\mathrm{U}(1)$. Theorem~\ref{thm:adams-unique} shows that
no other dimension in the cascade tower below $d_1$ could play this role.
Together, they fix $\mathrm{SU}(3)\times\mathrm{SU}(2)\times\mathrm{U}(1)$ as the unique
gauge group consistent with the cascade's topology at the gauge window.
\end{remark}

\begin{remark}[Hypercharge $\mathrm{U}(1)$ vs Paper~II's projective $\mathrm{U}(1)$]\label{rem:u1-vs-projective}
The cascade contains two structurally distinct $\mathrm{U}(1)$ groups,
both forced from cascade structure. They are abstractly isomorphic as
Lie groups (any $\mathrm{U}(1)$ is $S^1$) but physically distinct in
the cascade.
\begin{itemize}
\item \emph{Hypercharge $\mathrm{U}(1)$ (this paper, Theorem~\ref{thm:adams}
at $d=14$).} The unique nowhere-zero tangent vector field on $S^{13}$
($\rho(14)-1 = 1$). Lives at the specific cascade layer $d=14$. Acts
\emph{differentially} on charged states: a state with hypercharge $Y$
gets phase $e^{iY\phi}$. Observable in interference (shifts relative
phase by $(Y_2-Y_1)\phi$ between components of different hypercharge).
Becomes the photon after electroweak symmetry breaking via the
$\mathrm{SU}(2)$ obstruction at $d=13$.
\item \emph{Projective $\mathrm{U}(1)$ (Paper~II, $\S$7.5,
$\mathit{thm{:}U1{\text{-}}gauge{\text{-}}forced}$).} Generated by the
cascade's complex structure $J$ on the entire QM state space; closes
the projective state space to $\mathbb{C}P^{n-1}$ via the Hopf bundle.
Lives at \emph{every} cascade layer ($J$ is global). Acts
\emph{uniformly}: every state gets the same phase $e^{i\phi}$.
Gauge-trivial: preserves all relative phases; not observable in
interference. Forced under the cascade hypothesis by per-step propagator
consistency.
\end{itemize}
The two $\mathrm{U}(1)$s share a common \emph{generator} but differ in
their \emph{action}: distinct domains (charged-state subset at $d=14$
vs.\ full state space), distinct phase grading ($Y$-graded vs.\
uniform), and distinct observability (gauge boson vs.\ gauge-trivial).
The unification at the generator level is the content of the
structural connection below: both $\mathrm{U}(1)$s are 1-parameter
subgroups generated by the same cascade complex structure $J$, viewed
globally on $\mathbb{R}^{2n}$ (projective) vs.\ restricted to the
$d=14$ sphere $S^{13}$ (hypercharge), with the SM hypercharge weights
supplied separately by the Bott-period representation structure of
\S\S2.1--2.3 (chiral windows, complex layers, generator counting).
Numerical illustration:
\texttt{tools/research/projective\_u1\_vs\_hypercharge\_u1.py}.

\medskip
\noindent\emph{Structural connection: gauge-group identity closed,
eigenvalue conjecture falsified.} The hypercharge generator at $d=14$
\emph{is} the restriction of the global projective $J$ to the $d=14$
cascade sphere $S^{13}$, in the following precise sense.
\begin{enumerate}
\item \emph{Gauge group identity (closed).} $J$ on $\mathbb{R}^{2n}$
is antisymmetric (so $\langle Jx,x\rangle = 0$) and squares to
$-\mathrm{Id}$ (so $|Jx|^2 = -\langle x, J^2 x\rangle = 1$); this is
the content of Paper~II~\cite{paperII}'s
$\mathit{thm{:}complex}$ ($J$ from forced precession,
$J^2=-\mathrm{Id}$). Hence
$J|_{S^{2n-1}}$ is automatically a unit-norm tangent vector field at
every point. At $n=7$ (the $d=14$ layer), Theorem~\ref{thm:adams}
gives $\rho(14)-1 = 1$: any nowhere-zero tangent vector field on
$S^{13}$ spans the unique 1-dimensional Adams parallelism. Therefore
$J|_{S^{13}}$ \emph{is} the Adams field on $S^{13}$ (up to a constant
scalar), and the hypercharge $\mathrm{U}(1)$
$\exp(\theta J|_{S^{13}})$ and the projective $\mathrm{U}(1)$
$\exp(\theta J)$ are 1-parameter subgroups generated by the same $J$,
one restricted, one global.
\item \emph{Eigenvalue conjecture (falsified) and weight-spectrum gap.}
A stronger reading
---that $J|_{d=14}$'s eigenvalues directly reproduce the SM
hypercharge spectrum $Y\in\{-1,-\tfrac12,-\tfrac13,+\tfrac16,
+\tfrac23,+1\}$---is wrong: $J$ has only the eigenvalues $\pm i$, each
with multiplicity $n=7$, so $\exp(\theta J)$ acts as a uniform
$\pm$-phase rather than a $Y$-graded phase. The SM hypercharge
weights are therefore \emph{not} supplied by $J$ alone. They are also
not derived in \S\S2.1--2.3 of this paper: those subsections derive
the gauge factors at $d\in\{12,13,14\}$ (Theorem~\ref{thm:adams}),
the generator counts ($8+3+1$), $N_c=3$, and gauge-window uniqueness
(Theorem~\ref{thm:adams-unique}); they do not derive the SM Y
spectrum. This is the matter-rep / hypercharge gap recorded in
CLAUDE.md ``Known Quantitative Issues''. Closure status at $d=14$ is
the partial closure of Remark~\ref{rem:fund-or-trivial}: the
fundamental representation \emph{dimension} $\dim_{\mathbb{C}}=1$ is
forced; the specific $Y$ \emph{values} require a separate
cascade-internal derivation. A candidate route (cascade-internal
anomaly cancellation) is sketched in
\texttt{tools/research/cascade\_anomaly\_framework.py} and recorded
as an explicit open question in
Remark~\ref{rem:y-spectrum-open} below.
\item \emph{Coupling relation (open).} The cascade-internal coupling
at the $d=14$ layer is $\alpha(14)=R(14)^2/4\approx 0.0345$, used
downstream in Paper~IVb~\cite{paperIVb} for the universal shift
$\delta\Phi_{\mathrm{U}(1)}=\alpha(14)/\chi$. Its relation to the SM
running coupling $g_Y$ at any specific renormalisation scale is not
closed by the $J$-restriction argument and remains a separate open
question.
\end{enumerate}
No Part~IVa or Part~IVb prediction depends on item~(3); items~(1)
and~(2) tighten the structural picture from a speculative conjecture
to a closed identity (gauge-group) plus a falsified eigenvalue claim
plus an explicitly-stated weight-spectrum gap. Numerical
verification: \texttt{tools/research/projective\_u1\_vs\_hypercharge\_u1.py};
anomaly-framework scaffolding:
\texttt{tools/research/cascade\_anomaly\_framework.py};
open question: Remark~\ref{rem:y-spectrum-open}.
\end{remark}

\begin{remark}[Open question: cascade-internal derivation of the $Y$ spectrum]\label{rem:y-spectrum-open}
Remark~\ref{rem:fund-or-trivial} closes the fundamental representation
\emph{dimension} at $d=14$ ($\dim_{\mathbb{C}}=1$, forced by Adams'
theorem and the cascade complex structure $J$). It does not close the
specific $Y$ \emph{values} of the SM hypercharge spectrum
\[
Y\;\in\;\Bigl\{-1,\;-\tfrac{1}{2},\;-\tfrac{1}{3},\;
+\tfrac{1}{6},\;+\tfrac{2}{3},\;+1\Bigr\}.
\]
$J$'s eigenvalues are $\pm i$ (Remark~\ref{rem:u1-vs-projective}
item~(2)), so the fractional Y values cannot follow from the
$J$-restriction alone. This remark records the gap as an explicit
open question and sketches a candidate cascade-native route.

\medskip\noindent\textbf{The structural input the cascade has.}
\begin{itemize}
\item Per-particle $(V_{12},V_{13})$ assignments, fundamental or
  trivial of the gauge group at each layer
  (Remark~\ref{rem:fund-or-trivial}, Remark~\ref{rem:single-h-factor}).
\item Path-tensor structure
  $\mathcal{H}_{\text{matter}}=V_{12}\otimes V_{13}\otimes V_{14}$
  per particle (Remark~\ref{rem:path-tensor}).
\item Single-particle descent paths
  (Theorem~\ref{thm:forced-paths}).
\item The U(1) gauge structure at $d=14$ via $J$
  (Remark~\ref{rem:u1-vs-projective}).
\item Per-generation matter content of 15 Weyl fermions in the
  combinations $\{(\mathbf{3},\mathbf{2}),(\mathbf{3},\mathbf{1}),
  (\mathbf{3},\mathbf{1}),(\mathbf{1},\mathbf{2}),(\mathbf{1},
  \mathbf{1})\}$, each carrying a distinct $V_{14}$ that the cascade
  must specify.
\end{itemize}

\medskip\noindent\textbf{The candidate route: cascade-internal anomaly
cancellation.} The SM derives the $Y$ values up to overall
normalisation from four anomaly-cancellation conditions on
left-handed Weyl fermion content. A cascade-native analogue would
require each of the four to descend from a cascade structural
constraint, not from an imported QFT triangle diagram. The four
conditions and their cascade-native reading are:
\begin{enumerate}
\item[(I)] \emph{Gravitational $\times Y$:} $\sum_i n_i Y_i = 0$,
where $i$ runs over Weyl fermions per generation and $n_i$ is the
$\mathrm{SU}(3)\times\mathrm{SU}(2)$ rep dimension. Cascade-native
candidate: trace of the $V_{14}$-eigenvalue distribution over the
path-tensor matter content vanishes by descent-integral
well-definedness (the global cascade integral of a multi-particle
matter state at $d=14$ must be finite).
\item[(II)] \emph{$\mathrm{SU}(3)^2\times Y$:}
$\sum_{i\in\text{colored}} n_i^{\mathrm{SU}(2)} Y_i = 0$. Cascade-native
candidate: descent through the $d=12$ SU(3) layer requires the
sum of $V_{14}$ weights over color-fundamental matter to vanish, by
the analog of Theorem~\ref{thm:lefschetz}'s consistency of group
actions on $S^{11}$ for the multi-particle descent.
\item[(III)] \emph{$\mathrm{SU}(2)^2\times Y$:}
$\sum_{i\in\text{doublets}} n_i^{\mathrm{SU}(3)} Y_i = 0$.
Cascade-native candidate: parallel to (II) but at the $d=13$ layer,
through the Lefschetz obstruction on $S^{12}$.
\item[(IV)] \emph{$Y^3$:} $\sum_i n_i Y_i^3 = 0$. Cascade-native
candidate: third-order moment of the $V_{14}$ distribution under
the cubic structure of cascade descent integrals (open --- the
cascade has not yet identified a $Y^3$-like cubic structural
constraint).
\end{enumerate}
Each cascade-native candidate above is a structural conjecture, not
a theorem.

\medskip\noindent\textbf{Open question.} Identify cascade-internal
constraints (I')--(IV') such that:
\begin{itemize}
\item Each (X') is forced by cascade structure (descent paths,
  Lefschetz / hairy-ball obstructions, $J$-action on $\mathbb{C}^7$,
  or path-tensor compatibility), independently of any QFT input.
\item The system $\{(\text{I'}),(\text{II'}),(\text{III'}),
  (\text{IV'})\}$ together with one electric-charge anchor (e.g.\
  $Q_e=-1$ or a cascade-native equivalent) has a unique solution
  for the per-particle $V_{14}$ weights, matching the SM $Y$
  spectrum.
\end{itemize}
A first attempt at the framework is recorded in
\texttt{tools/research/cascade\_anomaly\_framework.py}, which sets
up the four anomaly conditions in cascade-native form and tests
which subsets of cascade primitives are sufficient to constrain
the $V_{14}$ weights. The script does not close the question; it
documents which of (I)--(IV) admit a cascade-native reading from
existing primitives and which require new structure.

\medskip\noindent\textbf{Structural prerequisite identified.}
A focused attempt to derive (I) ($\sum_i n_i Y_i = 0$) from existing
cascade primitives (descent-integral well-definedness, charge
conservation, the $J$-trace identity, Bott periodicity to spacetime,
and path-tensor consistency) was made in
\texttt{tools/research/cascade\_anomaly\_framework.py}, Step~2.5.
All five candidate routes fail. The shared failure mode is that
each anomaly condition (I)--(IV) requires a sum over multi-particle
matter content weighted by a specific gauge-anomaly weight function
(such as $n_i Y_i$ or $n_i Y_i^3$), and the cascade's existing
multi-particle structure (path-tensor of
Remark~\ref{rem:path-tensor} plus single-particle descent of
Theorem~\ref{thm:forced-paths}) does not yet supply such a weighted
sum. The structural ingredient that would supply it is a
gauge-coupled cascade fermion action --- the analogue of
$\bar\psi\gamma^\mu A_\mu\psi$ on the cascade lattice --- at the
gauge-window layers $d\in\{12,13,14\}$, and in particular at
$d=14$ for the four U(1)$_Y$ anomaly conditions.

The cascade currently has two relevant pieces of the fermion-action
structure:
\begin{itemize}
\item \emph{Largely closed:} the per-layer fermion partition function.
The Berezin construction with cascade-natural Dirac mass
$m(d)=\sqrt{\alpha(d)}$ gives $Z_f(d)=m(d)$ and
$Z_f(d)/Z_s(d)=1/(2\sqrt{\pi})$ exactly at every Dirac layer
(Part~IVb~\cite{paperIVb}, Remark on the Berezin partition-function
realisation, narrowing
Open Question~\texttt{oq:fermion-cascade-action}).
\item \emph{Substantially narrowed; one piece open:} the
gauge-coupling extension. Part~IVb~\cite{paperIVb}'s remark on the
proposed gauge-coupled action articulates the per-layer form
$g(d)\,\bar\psi T^a A^a \psi$ at gauge-window layers, with
$g(d) = \sqrt{\alpha(d)}$ matching the Yukawa $m(d) = \sqrt{\alpha(d)}$
($\sqrt{\alpha}$ universality), and \emph{derives} per-layer locality
of the cascade fermion from convergence of three cascade-source
readings. Per-layer Berezin reduction at $A=0$ is verified to machine
precision; the proposed gauge extension does not perturb the closed
mass-identification piece. The remaining open piece is reframed:
\emph{not} QFT triangle anomalies (which do not apply: per-layer
locality precludes derivative-coupled Dirac operators), but
\emph{path-tensor consistency} of the multi-layer gauge transformation
$U(1)_Y\times\mathrm{SU}(2)\times\mathrm{SU}(3)$ acting on the
single-particle path tensor $V_{12}\otimes V_{13}\otimes V_{14}$
(Remark~\ref{rem:path-tensor}). See
Part~IVb~\cite{paperIVb} Open Question~\texttt{oq:fermion-gauge-action}
and the proposal remark cited there for details.
\end{itemize}
The four anomaly conditions are therefore \emph{structurally
downstream} of Part~IVb's path-tensor-consistency reframing of
\texttt{oq:fermion-gauge-action} (the QFT triangle-anomaly framework
of the previous derivation attempt does not apply to per-layer-local
cascade fermions). The verifier
\texttt{tools/research/cascade\_anomaly\_framework.py} Step~2.5's
five-route audit remains valid as a NEGATIVE result on standard
cascade primitives; the path-tensor consistency framework is the
cascade-native replacement.

This narrows the Y-spectrum gap further: applying path-tensor
consistency to the cascade's EW-broken structure
(verifier
\texttt{tools/research/cascade\_path\_tensor\_consistency.py})
\emph{derives} the cascade-native relation
\[
Q \;=\; T_3 + Y
\]
at the unbroken U(1)$_{\rm em}$ at $d=4$ (the cascade-native
Gell-Mann--Nishijima identity, derived from the Higgs mechanism +
Part~IVb~\cite{paperIVb}~$\mathit{thm{:}weinberg}$). Combined with
the cascade-native pieces $N_c = 3$ at $d=12$, the SU(2) doublet
structure at $d=13$ ($Q$ differences within doublets equal $1$),
and $J$ eigenvalues $\pm 1$ at $d=14$ (U(1) fundamental unit $1$),
the cascade fixes: the denominator structure of $Q$ (and $Y$ in
$1/6$ units), the $Q$ differences within each multiplet, and the
relation $Q = T_3 + Y$. What remains undetermined cascade-natively
is the \emph{specific integer numerators} in SM electric charges
(e.g.\ $Q_u = +2/3$ as opposed to $-2/3$ or $0$). The Y-spectrum
gap is now reduced to one structurally-specific question: cascade-
native derivation of the integer charge numerators (equivalently,
of $Y_H$ at $d=13$, or of path-tensor trace identities on
$V_{12}\otimes V_{13}\otimes V_{14}$, or of the $J\,\pm 1$ eigenvalue
distribution over chirality-factorised matter content). See
Part~IVb~\cite{paperIVb} Open Question~\texttt{oq:fermion-gauge-action}
for the three attack routes articulated cascade-natively.

\medskip\noindent\emph{Status.} Open. Closing this question is
required to upgrade the d=14 entry of
Remark~\ref{rem:fund-or-trivial} from ``partially closed
(dimension)'' to ``closed (dimension and weights)''. Until then,
the cascade's claim of zero free parameters is precise about
predictions but not about the per-particle $V_{14}$ weight assignment.
\end{remark}

\begin{remark}[$\mathrm{SU}(3)$ algebra at $d_0=7$ via $G_2/\mathrm{SU}(3)$]\label{rem:su3-d7-algebra}
Theorem~\ref{thm:adams} derives the colour count $N_c=3$ from
$\rho(12)-1=3$ nowhere-zero tangent vector fields on $S^{11}$. A
direct bracket computation (verifier
\texttt{tools/\allowbreak verifiers/\allowbreak cascade\_d7\_su3\_bs\_closure.py}, Step~1)
shows these 3 fields, viewed as quaternionic right-multiplications
on $\mathbb{R}^{12}=\mathbb{H}^3$, close as the 3-dim Lie algebra
$\mathfrak{su}(2)\simeq\mathfrak{so}(3)$ under commutators
($[i,j]=2k$, $[j,k]=2i$, $[k,i]=2j$). This is 3-dimensional, not
the 8-dimensional $\mathfrak{su}(3)$ that the cascade's gauge group
identification requires.

The full $\mathfrak{su}(3)$ algebra is supplied cascade-internally
at a different distinguished layer: $d_0=7$. The cascade's
$S^6=\partial B^7$ is the unique higher-dimensional sphere
admitting an almost complex structure and is the homogeneous space
\[
S^6 \;=\; G_2/\mathrm{SU}(3),
\qquad \dim G_2 - \dim\mathrm{SU}(3) \;=\; 14-8 \;=\; 6 \;=\; \dim S^6.
\]
$G_2$ is the automorphism group of the octonions (Hurwitz at
$d=8$) and acts transitively on $S^6$ by isometries; $\mathrm{SU}(3)$
appears as the stabilizer of any unit imaginary octonion. Both
ingredients are cascade-forced: $d_0=7$ is the area maximum
(\cite{paperI}, Theorem~7.1; Paper~0 Theorem~5.2), and the
$G_2$ structure on $S^6$ is intrinsic to the cascade's $S^6$ at $d=7$.

Hence the cascade's $\mathrm{SU}(3)$ is distributed across two
distinguished layers, both cascade-forced:
\begin{enumerate}
\item \emph{Algebra} (8 generators) at $d_0=7$ via the $G_2/\mathrm{SU}(3)$
  homogeneous-space structure on $S^6$.
\item \emph{Fundamental representation dimension} ($N_c=3$) at $d=12$
  via Adams' theorem.
\end{enumerate}
Both layers are required; both are cascade-forced. The bracket
deficit at $d=12$ alone (3 vs.\ 8 generators) is filled by the
$d_0=7$ algebra structure.

\emph{Cascade-internal connection between $d=7$ and $d=12$.}  The two
layers contribute structurally different ingredients with consistent
dimensional content.  $d=12$ is the Bott image of the quaternion
layer ($d=12 - d=4 = 8$, one Bott period~$8$); the 3 Adams vector
fields on $S^{11}$ are the quaternionic units $\{i, j, k\}$ acting
on $\mathbb{R}^{12}=\mathbb{H}^3$ by right-multiplication, giving
$N_c=3$ but only an $\mathfrak{su}(2)$-isomorphic algebra (bracket
calculation above).  $d_0=7$ inherits its 8-dim
$\mathfrak{su}(3)$ algebra octonion-internally from $G_2 =
\mathrm{Aut}(\mathbb{O})$ acting on $S^6$.  The two are consistent
under the $\mathrm{SU}(N)$ family relation
$\dim\mathfrak{su}(N) = N^2-1$: $N=3$ uniquely satisfies rep dim~3
(from $d=12$) and algebra dim~8 (from $d_0=7$).
Cascade descent for any $\mathrm{SU}(3)$-anchored observable
(path $d=5..12$, Theorem~\ref{thm:forced-paths}) traverses $d_0=7$
on its way from gauge layer to observer, threading both
contributions through the same cascade integration.

This does not change any prediction of Sections~2--4. It
strengthens the cascade-internal forcing of the
gauge-group-identification step and supplies the structural source
layer for the $\alpha(7)/\chi^k$ correction family of Part~IVb~\S2:
the $\mathrm{SU}(3)$ algebra layer is the natural source for
quark-flavour observables that thread the $\mathrm{SU}(3)$ structure
(Cabibbo angle $\theta_C$ via $-\alpha(7)/\chi^2$; $b/s$ via
$-\alpha(7)/\chi^4$, see Part~IVb).
\end{remark}

\begin{remark}[Matter content composition via path tensor product]\label{rem:path-tensor}
The 30 Dirac d.o.f.\ per generation used implicitly throughout this
section ($Q_L$: 6, $u_R, d_R$: 3 each, $L_L$: 2, $e_R$: 1, summing
to 15 Weyl $=$ 30 Dirac without right-handed neutrinos) is the
path-tensor-product Hilbert dimension of the cascade matter content
under the gauge-window paths.

Cascade-internal forcing: the slicing recurrence $V_{d+1}=V_d
\cdot\sqrt{\pi}\,R(d+1)$ composes content multiplicatively across
layers, and the Adams gauge groups at $d\in\{12,13,14\}$ are
independent structures at distinct distinguished dimensions. The
unique cascade-native composition that respects (i) the
multiplicative slicing recurrence and (ii) the layer-independence
of Adams gauge groups is the tensor product:
\[
\mathcal{H}_{\text{matter}}\;=\;\bigotimes_{d\in\{12,13,14\}}V_{d,R(d)},
\]
where $V_{d,R(d)}$ is the matter content's representation at gauge
layer $d$. Direct sum is excluded by (i), wedge/symmetric products
by (ii). The verifier
\texttt{tools/verifiers/\allowbreak cascade\_path\_\allowbreak tensor\_\allowbreak product.py}
confirms that the SM rep content
($Q_L: \mathbf{3}\otimes\mathbf{2}\otimes Y$, etc.) gives 15 Weyl
per generation exactly.
\end{remark}

\begin{remark}[Fundamental-or-trivial at gauge-window layers]\label{rem:fund-or-trivial}
At each gauge-window layer $d\in\{12,13,14\}$, every matter unit's
representation $V_d$ under the gauge group at $d$ is the trivial
($\mathbf{1}$) or the fundamental of that group. No higher tensor
representations occur. The principle is closed layer by layer through
three independent cascade-native arguments.

\medskip\noindent\textbf{$d=12$ (SU(3)).}
The cascade space at $d=12$ is $\mathbb{R}^{12}=\mathbb{H}^3$, the
quaternionic Bott class inherited from $d=4$ by one Bott period~$8$.
The 3 quaternionic factors of $\mathbb{H}^3$ are the $N_c=3$ colors of
SU(3)'s fundamental representation, supplied by Adams' nowhere-zero
tangent fields on $S^{11}$ as quaternionic right-mults
(Remark~\ref{rem:su3-d7-algebra}: $\rho(12)-1=3$ generates the
$\mathfrak{su}(2)$ bracket structure on $\mathbb{H}^3$; the full
$\mathfrak{su}(3)$ algebra is supplied at the distinguished area
maximum $d_0=7$ via $G_2/\mathrm{SU}(3)$). Matter at $d=12$ lives in
this cascade space, so its SU(3) representation is one of the
SU(3)-irreducible components of $\mathbb{H}^3$ as an SU(3)-module: the
color-fundamental ($V_{12}=\mathbf{3}$) or the color-singlet
($V_{12}=\mathbf{1}$, present in the SU(3)-trivial sector). Closed
structurally; verifier
\texttt{tools/research/cascade\_fundamental\_rep\_lemma.py}, $d=12$
case.

\medskip\noindent\textbf{$d=13$ (SU(2)).}
Closed by Remark~\ref{rem:single-h-factor} below. Cascade-gauged
SU(2) at $d=13$ is the $d=12$ right-mult algebra $\{R_i,R_j,R_k\}$
extended trivially to the slicing axis (Adams' $\rho(13)-1=0$ admits
no native source). Two cascade-native routes restrict $V_{13}\in
\{\mathbf{1},\mathbf{2}\}$: color-charged matter
($V_{12}=\mathbf{3}$) via the Spin$(12)$ Dirac decomposition on
$\mathbb{H}^3$ under diagonal SU(2)$_R$; color-singlet matter
($V_{12}=\mathbf{1}$) via path-tensor layer independence
(Remark~\ref{rem:path-tensor}) combined with single-particle
descent (Theorem~\ref{thm:forced-paths}). The Spin$(12)$ Dirac route
is the explicit instance; the layer-independent route is the principle
applied at $d=13$ alone.

\medskip\noindent\textbf{$d=14$ (U(1)).}
Partially closed. The cascade complex structure $J$ on
$\mathbb{R}^{14}$ (with $J^2=-\mathrm{Id}$, supplied by Paper~II's
precession-induced complex structure) gives $\mathbb{R}^{14}=
\mathbb{C}^7$. Adams' theorem on $S^{13}$ yields $\rho(14)-1=1$, the
unique nowhere-zero tangent field; this field is $J|_{S^{13}}$ up to
scale. U(1) acts on each complex direction of $\mathbb{C}^7$ as
$\exp(\theta J)$ with eigenvalues $\pm i$, generating charges $\pm 1$.
The fundamental representation dimension $\dim_{\mathbb{C}}=1$ is
therefore forced. \emph{Open piece:} the specific SM hypercharge
values $Y\in\{-1,-\tfrac{1}{2},-\tfrac{1}{3},+\tfrac{1}{6},
+\tfrac{2}{3},+1\}$ are NOT forced by this argument
($J$'s eigenvalues are $\pm 1$, not fractional). Closing the
$Y$-spectrum cascade-internally requires a separate derivation
(e.g.\ cascade anomaly cancellation, or embedding into a richer
representation that picks specific $Y$ values); see CLAUDE.md
``Known Quantitative Issues'', matter-rep / hypercharge gap.

\medskip\noindent\emph{Status and outlook.}
The three layer arguments are independent: they invoke different
cascade structures (Bott quaternionic class at $d=12$, right-mult
extension at $d=13$, cascade complex structure at $d=14$). A
first-principles unification --- a single cascade theorem yielding
fundamental-or-trivial at all three gauge-window layers --- remains
open. The present remark aggregates the three closures under a single
referenceable name. Verifier:
\texttt{tools/research/cascade\_fundamental\_rep\_lemma.py}
(per-layer status; closes $d=12$ structurally and $d=14$ partially;
$d=13$ now closed by Remark~\ref{rem:single-h-factor}).
\end{remark}

\begin{remark}[Single-$\mathbb{H}$-factor matter forces $V_{13}\in\{\mathbf{1},\mathbf{2}\}$ at $d=13$]\label{rem:single-h-factor}
Remark~\ref{rem:path-tensor} composes matter content as a tensor product
$V_{12}\otimes V_{13}\otimes V_{14}$ across the gauge window, with one
tensor factor per gauge layer per particle. Combined with the
descent-path singularity of Section~\ref{sec:forced-paths}
(Theorem~\ref{thm:forced-paths}: each cascade observable threads exactly
one path), this forces a structural restriction at $d=13$: the matter
unit's $\mathrm{SU}(2)$ representation $V_{13}$ is either trivial
($\mathbf{1}$) or fundamental ($\mathbf{2}$), never a higher rep.

\medskip\noindent\textbf{Cascade-gauged $\mathrm{SU}(2)$ at $d=13$.}
Adams' theorem at $d=13$ gives $\rho(13)-1=0$, so no nowhere-zero
tangent field on $S^{12}$ sources the $\mathrm{SU}(2)$ generators
directly at $d=13$. The 3 generators of $\mathrm{SU}(2)$ at $d=13$ are
the $d=12$ right-multiplication algebra $\{R_i,R_j,R_k\}$, extended
trivially to the slicing axis at $d=13$:
\begin{itemize}
\item the $d=12$ right-mults already close as $\mathfrak{su}(2)$
(Remark~\ref{rem:su3-d7-algebra}: $[i,j]=2k$, $[j,k]=2i$, $[k,i]=2j$),
with $\dim\mathfrak{su}(2)=3$ matching the number of $\mathrm{SU}(2)$
generators required;
\item the right-mult action on $\mathbb{R}^{12}\subset\mathbb{R}^{13}$
vanishes on the two-point intersection $S^{12}\cap(\text{slicing axis})$,
consistent with the Lefschetz obstruction at $d=13$
(Theorem~\ref{thm:lefschetz}).
\end{itemize}

\medskip\noindent\textbf{Underlying principle.}
The rep-dimension restriction $V_{13}\in\{\mathbf{1},\mathbf{2}\}$ is
the $d=13$ instance of the principle stated in
Remark~\ref{rem:fund-or-trivial}: matter at each gauge-window layer
transforms in the fundamental or trivial representation of the gauge
group at that layer. The fundamental of SU(2) is the doublet
($\mathbf{2}$); the trivial is the singlet ($\mathbf{1}$). The
present remark closes the $d=13$ instance through two parallel
cascade-internal routes, which together establish the principle for
both SU(3) sectors of matter: color-charged matter
($V_{12}=\mathbf{3}$) via the Spin$(12)$ Dirac structure on
$\mathbb{H}^3$ (Route (i) below); color-singlet matter
($V_{12}=\mathbf{1}$) via path-tensor layer independence
(Remark~\ref{rem:path-tensor}) and single-particle descent
(Theorem~\ref{thm:forced-paths}) (Route (ii) below).

\medskip\noindent\textbf{Route (i): color-charged matter ($V_{12}=\mathbf{3}$).}
A matter unit transforming in the SU(3) fundamental occupies one
specific $\mathbb{H}$-factor of $\mathbb{H}^3$ (its color, per
Remark~\ref{rem:su3-d7-algebra}). The Spin$(12)$ Dirac on
$\mathbb{R}^{12}=\mathbb{H}^3$ factorises as $\mathrm{Spin}(4)^{\otimes 3}$
across the three $\mathbb{H}$-factors, with each Spin$(4)$ Dirac
splitting into Weyl$_+$ (SU(2)$_R$-singlet) and Weyl$_-$
(SU(2)$_R$-doublet) under the diagonal SU(2)$_R$
(verifier \texttt{tools/research/cascade\_route\_c\_d13\_dirac.py}).
Because the matter unit is concentrated in a single
$\mathbb{H}$-factor, it picks up exactly one Spin$(4)$-Weyl piece at
the corresponding tensor slot, and the other two Spin$(4)$ slots are
inert from the matter-unit perspective. Hence:
\begin{itemize}
\item zero Weyl$_-$ factors $\Rightarrow$ $V_{13}=\mathbf{1}$ (singlet);
\item one Weyl$_-$ factor $\Rightarrow$ $V_{13}=\mathbf{2}$ (fundamental).
\end{itemize}
The triplet $\mathbf{3}$ requires two Weyl$_-$ factors via
Clebsch--Gordan $\tfrac{1}{2}\otimes\tfrac{1}{2}=0\oplus 1$; the
quartet $\mathbf{4}$ requires three via
$\tfrac{1}{2}\otimes\tfrac{1}{2}\otimes\tfrac{1}{2}=\tfrac{1}{2}\oplus
\tfrac{1}{2}\oplus\tfrac{3}{2}$. Both configurations would require the
matter unit to occupy multiple $\mathbb{H}$-factors of $\mathbb{H}^3$
simultaneously, contradicting the SU(3)-fundamental concentration.
Triplet and quartet are excluded.

\medskip\noindent\textbf{Route (ii): color-singlet matter ($V_{12}=\mathbf{1}$).}
A matter unit trivial under SU(3) is not concentrated in any
$\mathbb{H}$-factor of $\mathbb{H}^3$; the Spin$(12)$ Dirac analysis of
Route (i) does not directly govern its SU(2) rep. The cascade-internal
argument at $d=13$ proceeds independently of the $d=12$ $\mathbb{H}^3$
structure:
\begin{itemize}
\item By the path-tensor decomposition of
Remark~\ref{rem:path-tensor},
$\mathcal{H}_{\text{matter}}=V_{12}\otimes V_{13}\otimes V_{14}$, with
factors at distinct distinguished cascade layers. $V_{12}$ and
$V_{13}$ are independent: $V_{12}=\mathbf{1}$ does not constrain
$V_{13}$ to any specific value, only the layer-$d=13$ gauge structure does.
\item The SU(2) gauge structure at $d=13$ is the
right-multiplication algebra extended trivially to the slicing axis
(paragraph above). Its irreducible representations are the
fundamental ($\mathbf{2}$) and the trivial ($\mathbf{1}$), plus higher
tensor reps ($\mathbf{3},\mathbf{4},\dots$) constructed multiplicatively.
\item Applying the fundamental-or-trivial principle at $d=13$ alone
restricts $V_{13}\in\{\mathbf{1},\mathbf{2}\}$. Higher reps would
require multi-particle tensor structure at the layer, which the
single-particle cascade descent (Theorem~\ref{thm:forced-paths}: each
matter unit threads exactly one path) excludes.
\end{itemize}
For color-singlet matter the triplet/quartet exclusion is therefore
\emph{direct from the fundamental-or-trivial principle plus
single-path descent}, not from the H$^3$-Dirac Clebsch--Gordan analysis
of Route (i).

\medskip\noindent\emph{Why the two routes coincide.} Both routes
restrict $V_{13}\in\{\mathbf{1},\mathbf{2}\}$ but through different
cascade structures: Route (i) uses the explicit
$\mathrm{Spin}(4)^{\otimes 3}$ tensor decomposition of the Dirac on
$\mathbb{H}^3$; Route (ii) uses the path-tensor layer independence
plus the fundamental-or-trivial principle at the gauge layer alone.
The agreement is a consistency check on the principle: where Route (i)
applies (color-charged matter), the H$^3$-Dirac structure realises the
fundamental-or-trivial restriction explicitly; where it does not
(color-singlet matter), the same restriction follows from the
layer-independent statement of the principle. The cascade-internal
content of the two routes is the same; only the path of derivation
differs.

\medskip\noindent\textbf{Connection to SM matter content.}
The Spin$(12)$ Dirac (built as Spin$(4)^{\otimes 3}$ on
$\mathbb{R}^{12}=\mathbb{H}^3$, dim $64$ complex) decomposes under
diagonal $\mathrm{SU}(2)_R$ chirality-by-chirality:
\begin{center}
\begin{tabular}{cccc}
\toprule
& Singlets & Doublets & Triplets / Quartet \\
\midrule
Spin$(12)$-Weyl$_+$ (chirality $+1$): & 14 & --- & 6 triplets \\
Spin$(12)$-Weyl$_-$ (chirality $-1$): & --- & 14 & 1 quartet \\
\bottomrule
\end{tabular}
\end{center}
This is the explicit decomposition relevant to Route~(i): on the
color-charged Spin$(12)$ Dirac sector, the
fundamental-or-trivial restriction selects the
$\{\text{singlet},\text{doublet}\}$ subset (i.e.\ $n_S\in\{0,1\}$
tensor combinations), excluding the $6$ triplets and the $1$ quartet.
Combined with the chirality split, this reproduces the SM
SU$(2)_L$ parity-violation pattern: $L$-handed states are doublets
($V_{13}=\mathbf{2}$), $R$-handed are singlets ($V_{13}=\mathbf{1}$).
For Route~(ii) matter (color-singlets), there is no Spin$(12)$
Dirac realisation to tabulate; the same $V_{13}\in\{\mathbf{1},
\mathbf{2}\}$ restriction applies via the principle directly. Either
way, triplet and quartet representations are excluded from the SM
matter spectrum cascade-internally. Numerical verifiers:
\texttt{tools/research/cascade\_fundamental\_rep\_lemma.py}
(per-layer status of the unified principle),
\texttt{tools/research/cascade\_route\_c\_d13\_dirac.py} (full Spin$(12)$
Dirac decomposition under diagonal SU(2)$_R$), and
\texttt{tools/research/cascade\_rep\_selection\_gap1.py} (decomposition
by number of Weyl$_-$ tensor factors and the Route (i) selection).

\medskip\noindent\emph{What this closes and what remains.}
The remark closes the rep-dimension question at $d=13$
($V_{13}\in\{\mathbf{1},\mathbf{2}\}$, no higher reps) for both SU(3)
sectors of matter. It does not address (a)~the specific multiplicities
of $\mathrm{SU}(2)$ singlets and doublets that match SM matter content
per generation (4 doublets $+$ 7 singlets per generation, vs.\ the
cascade's 14 of each in the Route (i) sector), and (b)~the specific SM
hypercharge values at $d=14$. Both remain open structural questions;
see CLAUDE.md ``Known Quantitative Issues'' for the broader
matter-rep gap. The fundamental-or-trivial principle itself is closed
at $d=12$ and $d=13$, partially closed at $d=14$ (dimension forced,
charge spectrum open); a tight first-principles proof of the principle
across all three gauge layers from a single cascade theorem is a
natural follow-up direction.
\end{remark}

\section{Symmetry Breaking from the Hairy Ball Theorem}
\label{sec:breaking}

\subsection{Even and odd spheres in the cascade}

The hairy ball theorem (Poincar\'{e}, 1885~\cite{poincare}): every continuous tangent vector
field on an even-dimensional sphere $S^{2n}$ has at least one zero. Odd-dimensional spheres
$S^{2n+1}$ admit nonvanishing vector fields. The Euler characteristic
$\chi(S^{2n})=2$, $\chi(S^{2n+1})=0$.

The cascade at dimension $d$ operates on $S^{d-1}$. The slicing direction defines a tangent
vector field on this sphere. Therefore:
\begin{itemize}
\item $d$ odd $\Rightarrow$ $S^{d-1}$ even-dimensional $\Rightarrow$ tangent field has
forced zero.
\item $d$ even $\Rightarrow$ $S^{d-1}$ odd-dimensional $\Rightarrow$ tangent field can be
nonvanishing.
\end{itemize}

\begin{center}
\begin{tabular}{ccccc}
\toprule
$d$ & Sphere & Parity & Hairy ball & Gauge/Status \\
\midrule
12 & $S^{11}$ & odd  & smooth & $\mathrm{SU}(3)$ unbroken \\
13 & $S^{12}$ & even & ZERO   & $\mathrm{SU}(2)$ broken \\
14 & $S^{13}$ & odd  & smooth & $\mathrm{U}(1)$ unbroken \\
\bottomrule
\end{tabular}
\end{center}

\subsection{The Lefschetz obstruction: a group-theoretic theorem}

The hairy ball theorem shows that any tangent vector field on $S^{12}$ has a zero. A
stronger statement is available: no connected compact Lie group can act freely on any
even-dimensional sphere. This upgrades the breaking from a field-theoretic observation to
a group-theoretic theorem.

\begin{theorem}[Lefschetz obstruction on even spheres]\label{thm:lefschetz}
Let $f\colon S^{2n}\to S^{2n}$ be any continuous orientation-preserving map. Then
the Lefschetz number satisfies $L(f)=2\neq 0$, and $f$ has at least one fixed point.
\end{theorem}
\begin{proof}
The Lefschetz number is
\[
L(f)=\sum_{k=0}^{2n}(-1)^k\,\mathrm{tr}\!\bigl(f_*\mid_{H^k(S^{2n};\mathbb{Q})}\bigr).
\]
The rational cohomology of $S^{2n}$ is $H^0=H^{2n}=\mathbb{Q}$ and $H^k=0$ otherwise.
The map $f$ acts as the identity on $H^0$ and as multiplication by $\deg(f)$ on $H^{2n}$.
Since $f$ is orientation-preserving, $\deg(f)=+1$. Therefore
\[
L(f)=(-1)^0\cdot 1+(-1)^{2n}\cdot 1=1+1=2\neq 0.
\]
By the Lefschetz fixed-point theorem~\cite{lefschetz}, $L(f)\neq 0$ implies $f$ has at
least one fixed point.
\end{proof}

\begin{corollary}[No free Lie group actions on even spheres]\label{cor:no-free}
Let $G$ be a connected compact Lie group acting on $S^{2n}$ by orientation-preserving
diffeomorphisms. Then every element $g\in G$ has at least one fixed point; in particular,
$G$ does not act freely on $S^{2n}$.
\end{corollary}
\begin{proof}
Since $G$ is connected, every $g\in G$ is homotopic to $\mathrm{id}$ via a path $g_t$.
Each $g_t$ is orientation-preserving (as $G$ is connected and $\mathrm{id}$ is
orientation-preserving), so $\deg(g)=+1$ for all $g\in G$. By
Theorem~\ref{thm:lefschetz}, $L(g)=2\neq 0$, so every $g$ has a fixed point.
\end{proof}

\begin{remark}
Corollary~\ref{cor:no-free} is stronger than the hairy ball theorem alone. The hairy ball
theorem shows that the \emph{slicing vector field} has a zero at $d=13$. The Lefschetz
corollary shows that \emph{no connected compact Lie group can act freely on} $S^{12}$,
regardless of the field configuration. SU(2) breaking is a group-theoretic obstruction,
not merely a field-theoretic observation. Both the original hairy ball argument and the
Lefschetz theorem give the same physical conclusion; the Lefschetz version is the stronger
statement.
\end{remark}

\subsection{The gauge window parity pattern}

\begin{theorem}[Standard Model breaking pattern]\label{thm:breaking}
In the gauge window $\{12,13,14\}$: $\mathrm{SU}(3)\times\mathrm{U}(1)_{\rm em}$ is
unbroken and $\mathrm{SU}(2)_L$ is broken.
\end{theorem}
\begin{proof}
$d=12$ is even, so $S^{11}$ is odd-dimensional. The slicing vector field has no forced
zero; $\mathrm{SU}(3)$ is unbroken. $d=13$ is odd, so $S^{12}$ is even-dimensional.
By Corollary~\ref{cor:no-free}, no connected compact Lie group (in particular
$\mathrm{SU}(2)$) can act freely on $S^{12}$; $\mathrm{SU}(2)$ is broken. $d=14$ is
even, so $S^{13}$ is odd-dimensional; $\mathrm{U}(1)$ is unbroken. This is exactly the
Standard Model pattern: $\mathrm{SU}(3)_{\rm colour}\times\mathrm{U}(1)_{\rm em}$
unbroken, $\mathrm{SU}(2)_L$ broken.
\end{proof}

\subsection{The Higgs as the hairy ball zero}

The forced zero on $S^{12}$ at $d=13$ is the cascade's Higgs mechanism. The zero is a
topological obstruction, not a dynamical field acquiring a vacuum expectation value. Its
existence is guaranteed by $\chi(S^{12})=2$; its location on $S^{12}$ is determined by the
cascade's slicing geometry. By the Poincar\'{e}--Hopf theorem~\cite{hopf}, the sum of the
indices of all zeros of a tangent vector field on a compact manifold $M$ equals $\chi(M)$.
For $S^{12}$: $\chi(S^{12})=2$. The tangent field $v(\theta)=\sin\theta$ has zeros at
$\theta=0$ and $\theta=\pi$, each with index 1, summing to $\chi=2$.

\begin{remark}[Cascade vs Standard Model Higgs]
In the Standard Model, the Higgs field is a complex scalar doublet $\phi$ with potential
$V(\phi)=\lambda(|\phi|^2-v^2/2)^2$. The VEV is chosen by spontaneous symmetry breaking,
and $m_H=\sqrt{2\lambda}\,v$ depends on the quartic coupling $\lambda$, a free parameter.
In the cascade, the Higgs is the hairy ball zero---a topological obstruction whose existence
is a theorem, not a choice. The VEV is determined by the geometry of $S^{12}$ (the equator,
where $\sin\theta$ is maximal). The mass ratio $m_H/m_W=\pi/2$ is a geodesic distance, not
a coupling constant. The cascade replaces one free parameter ($\lambda$) with one geometric
fact ($\pi/2$).
\end{remark}

\subsection{The Higgs mass from the geodesic distance}

The $\mathrm{SU}(2)$ gauge field on $S^{12}$ is a tangent vector field $v(\theta)=\sin\theta$,
where $\theta$ is the geodesic distance from the forced zero at the pole. Two points are
geometrically distinguished:
\begin{itemize}
\item The zero at $\theta=0$ (north pole): $\sin 0=0$. The Higgs lives here.
\item The maximum at $\theta=\pi/2$ (equator): $\sin(\pi/2)=1$. The VEV lives here. This
is where $\mathrm{SU}(2)$ is maximally broken and the $W$ boson acquires mass.
\end{itemize}

\begin{theorem}[Higgs-to-$W$ mass ratio]\label{thm:higgs}
On $S^{12}$ at $d=13$, the mass ratio of the Higgs boson to the $W$ boson equals the
geodesic distance from the hairy ball zero to the VEV:
\[
m_H/m_W = \pi/2.
\]
\end{theorem}
\begin{proof}
The $W$ boson mass is set by the gauge field's value at the VEV:
$m_W\propto|v(\pi/2)|=\sin(\pi/2)=1$. This defines the energy scale $R$ of $\mathrm{SU}(2)$
breaking. The Higgs mass is set by the energy of the field configuration connecting the zero
to the VEV. On a sphere of radius $R$, the arc length from pole to equator is $R\times\pi/2$.
The ratio is purely angular: $m_H/m_W=\pi/2$.

Numerically: $m_H=m_W\times\pi/2=80.38\times 1.5708=126.3$ GeV. Observed: 125.25 GeV.
Deviation: 0.80\%.
\end{proof}

\begin{remark}[The quarter turn appears for the fourth time]
The angle $\pi/2$ appears in four roles: (i)~the slicing axis is perpendicular to the
equatorial hyperplane, forcing $\Gamma(\tfrac{1}{2})=\sqrt{\pi}$ in the slicing
recurrence~\cite{paperI}; (ii)~consecutive slicing axes are perpendicular, forcing the
precession $\alpha=\pi/2$~\cite{paperII}; (iii)~two quarter-turns give $J^2=-\mathrm{Id}$,
producing complex structure~\cite{paperII}; (iv)~the hairy ball zero and the VEV are
separated by $\pi/2$ on $S^{12}$, fixing $m_H/m_W$. All four are the same geometric
fact---orthogonality---applied to different objects in the cascade.
\end{remark}

\subsection{The cascade obstruction map}

Combining the Clifford classification, the hairy ball theorem, and the gauge window
structure, the full layer-by-layer topology from the observer to the third Bott window is:

\begin{center}
\resizebox{\textwidth}{!}{%
\begin{tabular}{cccccl}
\toprule
$d$ & $d\bmod 8$ & Type & Sphere & $\chi$ & Role \\
\midrule
4  & 4 & Weyl  & $S^3$  & 0 & Observer \\
5  & 5 & Dirac & $S^4$  & 2 & Gen 3 ($\tau,b,\nu_\tau$); zero \\
6  & 6 & Weyl  & $S^5$  & 0 & --- \\
7  & 7 & Real  & $S^6$  & 2 & $d_0$ (zero forced; real spinor, no complex amplitude) \\
8--11 & 0--3 & Real & $S^7$--$S^{10}$ & var & --- \\
12 & 4 & Weyl  & $S^{11}$ & 0 & $\mathrm{SU}(3)$ \\
13 & 5 & Dirac & $S^{12}$ & 2 & Gen 2 ($\mu,s,c$); $\mathrm{SU}(2)$ broken \\
14 & 6 & Weyl  & $S^{13}$ & 0 & $\mathrm{U}(1)$ \\
15--19 & 7--3 & Real & $S^{14}$--$S^{18}$ & var & $d_1$ at 19 \\
20 & 4 & Weyl  & $S^{19}$ & 0 & --- \\
21 & 5 & Dirac & $S^{20}$ & 2 & Gen 1 ($e,d,u$); zero \\
22 & 6 & Weyl  & $S^{21}$ & 0 & --- \\
\bottomrule
\end{tabular}%
}
\end{center}

Three structural features are visible:
\begin{enumerate}
\item Hairy ball zeros appear only at Dirac layers ($d\bmod 8=5$, which are odd, so $S^{d-1}$
is even). These are the generation layers and the topological obstructions of~\cite{paperIVb}.
\item The gauge window $\{12,13,14\}$ sits between the first ($d=5$) and second ($d=13$)
hairy ball zeros, centred on the self-dual crossing.
\item Real layers carry no zeros and no complex structure. They are transparent to fermion
propagation---contributing only geometric attenuation, not topological obstruction.
\end{enumerate}

\section{Three Generations from the Phase Transition}

\subsection{Generation candidates}

\begin{definition}[Generation layer]
A dimension $d$ is a generation layer if:
\begin{itemize}
\item[(G1)] $d\bmod 8=5$: the minimal spinor is complex Dirac, and
\item[(G2)] $d$ is odd: $S^{d-1}$ is even-dimensional, so the hairy ball theorem forces
a zero.
\end{itemize}
Since 5 is odd, all complex Dirac layers automatically satisfy (G2). The generation layers
are $d=5,13,21,29,37,\ldots$, spaced by one Bott period.
\end{definition}

\subsection{The \texorpdfstring{$d_1=19$}{d\_1=19} phase transition as cutoff}

The sphere-area decay rate $p(d)=\tfrac{1}{2}\psi((d+1)/2)-\tfrac{1}{2}\ln\pi$ crosses
the first threshold $c_1=\tfrac{1}{2}\ln\pi$ at the continuous value $d^*_1=19.731$.
Below $d_1$, the cascade is subcritical: each step loses less than one factor of $\sqrt{\pi}$.
Above $d_1$, it is supercritical: exponential suppression sets in.

\begin{center}
\begin{tabular}{cccl}
\toprule
Layer & $d$ & $p(d)-c_1$ & Regime \\
\midrule
Gen 3 & 5  & $-0.683$ & Subcritical ($\tau,b,\nu_\tau$) \\
Gen 2 & 13 & $-0.208$ & Subcritical ($c,s,\mu$) \\
Gen 1 & 21 & $+0.031$ & Supercritical ($u,d,e$) \\
Gen 0 & 29 & $+0.192$ & Supercritical (absent) \\
\bottomrule
\end{tabular}
\end{center}

\subsection{Why three and only three}

\begin{theorem}[Three charged-fermion generations, with a suppressed fourth Bott layer]\label{thm:generations}
The number of observable charged-fermion generations is exactly three
($d=5,13,21$). The fourth Bott fermion layer at $d=29$ exists as a
cascade primitive (Bott periodicity is 8, and $21+8=29$ is the next
Dirac layer) but is suppressed in charged-fermion mass spectra by a
factor of $\sim 289$ relative to Generation~1. Its residual role is
to source neutrino masses (Part~IVb, Roadmap~item~2: $m_{29}\times
\alpha(21)/\chi^8 = 0.0493$~eV for the heaviest neutrino, $-1\%$ on
$\Delta m^2$).
\end{theorem}
\begin{proof}
Generation 3 ($d=5$) and Generation 2 ($d=13$) are subcritical:
$p(d)<c_1$ at both layers, so their propagator amplitudes from $d=4$
are $O(1)$ and unsuppressed.

Generation 1 ($d=21$) is barely supercritical: $p(21)-c_1=+0.031$,
only 1.3 steps past the threshold in units of the transition width
$1/p'(d_1)\approx 39.5$. Its propagator amplitude is suppressed but
nonzero.

The fourth Bott layer at $d=29$ is 9.3 steps past threshold; its
charged-fermion propagator amplitude relative to Generation~1 is
suppressed by a factor of $\sim 289$, making it invisible in the
observed charged-fermion mass spectrum. The transition at $d_1=19$
---derived in~\cite{paperI} from the Gamma function alone---is the
wall that kills the fourth \emph{charged-fermion} generation while
leaving $d=29$ available as a cascade-primitive layer; Part~IVb
(Roadmap~item~2) uses this layer as the source of the heaviest
neutrino mass via the chirality-filter exponent $\chi^8=256$,
recovering $\sqrt{\Delta m^2_{\mathrm{atm}}} = 0.0495$~eV to $-1\%$.

The transition width deserves comment. The derivative
$p'(d)=\tfrac{1}{4}\psi^{(1)}((d+1)/2)>0$ is the trigamma function,
strictly positive everywhere. At $d_1=19$:
$p'(19)=\tfrac{1}{4}\psi^{(1)}(10)\approx 0.0253$. The transition
width $\Delta d\equiv 1/p'(d_1)\approx 39.5$ layers means the
transition is smooth, not sharp. Generation 1 at $d=21$ is only
$\Delta d_1=(21-19.73)/39.5=0.032$ transition widths past the
threshold. The fourth Bott layer at $d=29$ is 0.24 widths
past---modest, but the exponential nature of the supercritical
regime makes even this small overshoot devastating for
charged-fermion amplitudes:
$\exp(8\times(p(29)-c_1))=\exp(8\times 0.192)\approx 4.6$ additional
suppression per extra generation.
\end{proof}

\begin{remark}[Derivation status of the generation count]\label{rem:sp30-status}
Theorem~\ref{thm:generations} factorises into three statements with
three different derivation statuses:
\begin{enumerate}
\item \emph{Derived.} Exactly three observable charged-fermion
generations at $d\in\{5,13,21\}$: follows from the charged-fermion
amplitude suppression at $d=29$ (the $\sim 289$ factor from
$\exp(8\cdot(p(29)-c_1))$) and the phase-transition wall at
$d_1=19$. The three-versus-four \emph{charged} count is forced.
\item \emph{Derived.} The fourth Bott fermion layer at $d=29$ is a
cascade primitive: Bott periodicity is 8 and the Dirac layers are
$d\bmod 8 = 5$, so $\{5,13,21,29,\ldots\}$ is the cascade's
Dirac-layer sequence. The first three enter charged-fermion
channels; the fourth is suppressed there. No cascade theorem
eliminates $d=29$ as a layer.
\item \emph{Derived (Part~IVb).} The residual role of $d=29$ as a
neutrino-mass source: Part~IVb derives
$m_{29}\cdot\alpha(21)/\chi^8 = 0.0493$~eV for the heaviest
neutrino, matching $\sqrt{\Delta m^2_{\mathrm{atm}}} = 0.0495$~eV
to $-1\%$. The chirality-filter exponent $\chi^8 = 256$ is one full
Bott period of chirality filtering (cascade distance $29-21=8$).
\end{enumerate}
The earlier theorem statement ``the number of observable fermion
generations is exactly three'' and the earlier proof phrasing ``the
wall that kills the fourth generation'' implied that the $d=29$
layer is nonexistent. The three-step reading above keeps both
claims truthful: exactly three observable \emph{charged-fermion}
generations, and the $d=29$ layer partially realised as the
neutrino-mass source. An earlier passage in the abstract
(``$d=29$ is deep in the supercritical regime and absent'') is
correspondingly softened in the next commit.
\end{remark}

\section{Propagator Phases and the Quaternionic Structure}

The cascade propagator accumulates a forced phase of $\pi/2$ per step
(Theorem~6.1 of~\cite{paperII}). From $d=4$ to the three gauge layers:
\begin{align*}
d=12:&\quad 8 \text{ steps}\Rightarrow \text{phase}=4\pi \Rightarrow e^{i\cdot 4\pi}=+1,\\
d=13:&\quad 9 \text{ steps}\Rightarrow \text{phase}=9\pi/2 \Rightarrow e^{i\cdot 9\pi/2}=i,\\
d=14:&\quad 10 \text{ steps}\Rightarrow \text{phase}=5\pi \Rightarrow e^{i\cdot 5\pi}=-1.
\end{align*}
The three gauge layers carry the fourth roots of unity: $\{+1,i,-1\}$. Together with the
identity at $d=4$, the set $\{+1,i,-1,-i\}$ constitutes the unit quaternions $\pm 1,\pm i$
restricted to the complex plane. The missing element $k$ of the full quaternionic structure
$\mathbb{H}$ would require a fourth gauge factor, absent by the hairy ball theorem. The
Standard Model's gauge group is the maximal group consistent with the cascade's complex
(not quaternionic) structure.

\begin{remark}[Division algebra interpretation]
The cascade forces the quantum amplitude algebra to be $\mathbb{C}$ (real dimension~2), via
$J^2=-\mathrm{Id}$ from the forced precession~\cite{paperII}. The associative normed division
algebras over $\mathbb{R}$ are $\mathbb{R}$, $\mathbb{C}$, $\mathbb{H}$, $\mathbb{O}$
(Hurwitz theorem~\cite{hurwitz}; $\mathbb{O}$ is non-associative and excluded from quantum
mechanics by associativity of sequential measurements). The gauge window's quaternionic
phases $\{+1,i,-1\}$ are the $\mathbb{C}$-restriction of $\mathbb{H}$: the cascade reaches
toward $\mathbb{H}$ but cannot complete it because the hairy ball theorem allows only three
gauge factors. This is structurally parallel to Furey's observation that $\mathbb{C}\otimes
\mathbb{H}$ organises the electroweak sector: the cascade produces the same algebra from
its topology.
\end{remark}

\begin{remark}[The Gram identity for cascade layer overlaps]\label{rem:gram}
The $L^2$ overlap of cascade layer integrands $f_d(x)=(1-x^2)^{d/2}$ satisfies:
\[
\langle f_{d_1},f_{d_2}\rangle = \int_{-1}^1(1-x^2)^{(d_1+d_2)/2}dx = N(d_1+d_2).
\]
The Gram matrix of all cascade layer integrands is $G_{d_1 d_2}=N(d_1+d_2)$: the
cross-overlap of any two layers equals the cascade lapse evaluated at their summed dimension.
The overlap structure of cascade layers is therefore entirely encoded in the lapse function
itself; no new geometric objects enter.

Applied to the generation layers: $\hat{O}(5,13)=N(18)/\sqrt{N(10)N(26)}=0.948$,
$\hat{O}(5,21)=0.888$, $\hat{O}(13,21)=0.986$. The three generation wavefunctions are
nearly linearly dependent, spanning an effectively one-dimensional subspace consistent with
their common Dirac structure ($d\bmod 8=5$). The Gram determinant
$\det G_{\rm gen}\approx 0.002$ confirms near-singularity.
\end{remark}

\section{Forced Cascade Paths}\label{sec:forced-paths}

The preceding sections place specific dimensions on specific physical roles: the fermion generation layers at $d\in\{5,13,21\}$ (Theorem~\ref{thm:generations}), the gauge boson layers at $d\in\{12,13,14\}$ with $\mathrm{SU}(3)$ uniquely at $d=12$ (Theorem~\ref{thm:adams-unique}), and the observer host at $d_V=5$, the volume maximum of the cascade (\cite{paperI}, Theorem~7.1). The cascade paths used by Part~IVb~\cite{paperIVb} to compute mass spectra, gauge couplings, and precision predictions are not independent inputs: once these layers are fixed, every path is forced as the cascade descent between two already-determined endpoints. This section collects the statements.

\begin{definition}[Cascade potential]\label{def:cascade-potential}
The cascade potential at dimension $d\geq 5$ is
\[
\Phi(d) \;:=\; \sum_{d'=5}^{d} p(d'), \qquad p(d') = \tfrac{1}{2}\psi\!\Bigl(\tfrac{d'+1}{2}\Bigr) - \tfrac{1}{2}\ln\pi,
\]
the cumulative log decay rate of sphere areas from the volume maximum $d_V=5$ up to dimension $d$. The definition is normalised so that $\Phi(5)=p(5)$ is the initial slicing step out of the observer host.
\end{definition}

\begin{theorem}[Derived cascade paths]\label{thm:forced-paths}
Every cascade path invoked in Part~IVb~\cite{paperIVb} is forced by the layer assignments of this paper and Paper~I~\cite{paperI}. Specifically:
\begin{enumerate}
\item[(i)] \textbf{Mass ratios between fermion generations.} For two fermion generation layers $d_A<d_B$ with $d_A,d_B\in\{5,13,21\}$ (Theorem~\ref{thm:generations}), the geometric log mass ratio is the cascade potential difference between them:
\[
\log\!\left(\frac{m(d_A)}{m(d_B)}\right)_{\text{geometric}} \;=\; \Phi(d_B)-\Phi(d_A) \;=\; \sum_{d=d_A+1}^{d_B} p(d),
\]
where $m(d_A)$ is the heavier generation (smaller $d$, nearer the volume maximum, more cascade content) and $m(d_B)$ is the lighter. The sum runs over the $d_B-d_A$ slicing steps strictly between the two fermion layers. In particular, $m_\tau/m_\mu$ is forced onto $d=6..13$ and $m_\mu/m_e$ onto $d=14..21$. Both paths have $d_B-d_A=8$, the Bott period, because the generation layers are three consecutive points in the $d\bmod 8 = 5$ orbit. No other length is consistent with the generation placement.

\item[(ii)] \textbf{Gauge-anchored observables.} For any observable whose cascade home is a gauge boson layer $d_B\in\{12,13,14\}$ (Theorem~\ref{thm:adams}, Theorem~\ref{thm:adams-unique}), the geometric log attenuation from its gauge layer down to the observer at $d=4$ is the cascade potential at the gauge layer:
\[
\log\!\left(\frac{Q_{\text{obs}}}{Q_{\text{bare}}}\right) \;=\; \Phi(d_B) \;=\; \sum_{d=5}^{d_B} p(d),
\]
where the sum runs over all $d_B-4$ slicing steps from the volume maximum $d_V=5$ (observer host) up to the gauge boson's own layer $d_B$. For $\mathrm{SU}(3)$ at $d=12$ this is the path $d=5..12$, giving $\alpha_s(M_Z) = \alpha_{\mathrm{GUT}}\cdot\exp(\Phi(12))$ as used in Part~IVb~\cite{paperIVb} (strong coupling theorem). For $\mathrm{SU}(2)$ at $d=13$ this is $d=5..13$ and for $\mathrm{U}(1)$ at $d=14$ this is $d=5..14$, feeding the Weinberg angle calculation of \cite{paperIVb} (Weinberg angle theorem).

\item[(iii)] \textbf{Higgs-mediated scales.} The electroweak VEV $v$ and every fermion mass are built multiplicatively from $\alpha_s$ via the universal coupling $C=\alpha_s/(2\sqrt{\pi})$ (Part~IVb~\cite{paperIVb}, universal coupling theorem), so their cascade path is inherited from case~(ii) with $d_B=12$: the Higgs sits at the $\mathrm{SU}(3)$-adjacent hairy-ball zero on $S^{12}$ (Section~4.4 above), and the VEV propagates on $d=5..12$.

\item[(iv)] \textbf{Intra-generation ratios} (up- vs down-type quarks, or charged leptons vs neutrinos within a single generation) are determined by the chirality sub-structure at a single fermion layer, not by a cascade descent; they do not receive a path contribution.
\end{enumerate}
\end{theorem}

\begin{proof}
Case~(i). The generation layers $d=5,13,21$ are forced by (G1) complex Dirac minimal spinor at $d\bmod 8=5$ and (G2) odd $d$ ensuring an even-dimensional boundary sphere $S^{d-1}$, intersected with the $d_1=19$ phase transition cutoff (Theorem~\ref{thm:generations}). The geometric component of a mass ratio is the exponential of the cascade potential difference between the two layers. In Definition~\ref{def:cascade-potential}, $\Phi(d)$ accumulates $p(d')$ from $d'=5$; the difference $\Phi(d_B)-\Phi(d_A)$ for $d_A<d_B$ is therefore $\sum_{d=d_A+1}^{d_B} p(d)$, and no term at $d_A$ enters because it cancels between $\Phi(d_B)$ and $\Phi(d_A)$. Since the three generations are three consecutive points in the same mod-8 Bott orbit, the descent from one to the next always covers exactly $8$ layers. This fixes the path length without reference to the observer or any external scale.

Case~(ii). The gauge boson layers $d=12,13,14$ are forced by the Radon--Hurwitz numbers $\rho(d)-1$ on $S^{d-1}$ (Theorem~\ref{thm:adams}), and $d=12$ is the \emph{unique} dimension in $[5,d_1=19]$ realising $\rho(d)-1=3$ (Theorem~\ref{thm:adams-unique}). The observer host $d_V=5$ is the unique volume maximum of the cascade (\cite{paperI}, Theorem~7.1). A gauge-anchored observable lives at its gauge layer in the bulk and is observed at $d=4$; its cascade descent must therefore traverse every slicing step from the gauge layer $d_B$ down to the observer host $d_V=5$, contributing $p(d)$ at each. The resulting sum is $\Phi(d_B)=\sum_{d=5}^{d_B}p(d)$, of length $d_B-4$. No free parameters enter.

Case~(iii). Part~IVb's universal coupling theorem expresses every fermion mass and the electroweak VEV as products involving $\alpha_s$ and a topological factor. Since the topological factor is path-independent (Part~IVb~\cite{paperIVb}, obstruction primitive corollary: $2\sqrt{\pi}=N(0)\cdot\Gamma(\tfrac{1}{2})$), the path content of each such quantity is inherited from the cascade potential $\Phi(12)$ that builds $\alpha_s$.

Case~(iv) is not a descent statement: same-layer ratios depend on the spinor decomposition of the single sphere $S^{d-1}$ at that layer, not on a multi-layer cascade integral.
\end{proof}

\begin{remark}[What is not forced]\label{rem:not-forced}
Theorem~\ref{thm:forced-paths} derives the \emph{path} of every observable whose cascade home is a distinguished layer (generation or gauge boson). It does \emph{not} yet derive (a) the \emph{correction formula} to the leading-order $\exp(\Phi)$ on each path---this is the status of the Part~0 Supplement and the remaining observable-dependent question flagged in \cite{paperIVb}; (b) the absolute mass scale, which enters via the electroweak VEV and ultimately the reduced Planck mass as a dimensional anchor; (c) up-type quark mass ratios, which the paper classifies as Tier~4 pending a derivation of the chirality correction at the $\mathrm{SU}(3)$ layer. The theorem separates the path problem from the correction problem: with paths forced, any remaining precision deviation is isolated to the correction formula, not to the choice of cascade layers.
\end{remark}

\begin{remark}[Eight is Bott, not choice]\label{rem:bott-length}
The mass-ratio descent in case~(i) has length exactly $8$ (the Bott period of $O(d)$ stable homotopy): $d_B-d_A=8$ for both $(d_A,d_B)=(5,13)$ and $(13,21)$. This is not arithmetic coincidence: the generation layers are three consecutive points in the $d\bmod 8 = 5$ Bott orbit. The gauge-anchored descent in case~(ii) has length $d_B-4 \in\{8,9,10\}$ for $d_B\in\{12,13,14\}$---each gauge layer lies within one Bott window $\{d_V+7,d_V+8,d_V+9\}$ of the observer host $d_V=5$. Any observable at a more distant layer would lie beyond $d_1=19$ and be suppressed by the phase transition (\cite{paperI}, Theorem~8.4). The cascade's own internal coherence length is the Bott period, and every derived path respects it.
\end{remark}

\section{What This Paper Proves}

This paper proves, from the cascade's geometry and classical mathematical theorems:

\begin{enumerate}
\item The gauge group is $\mathrm{SU}(3)\times\mathrm{SU}(2)\times\mathrm{U}(1)$, derived
from the Bott mirror window $\{12,13,14\}$ and Adams' theorem.

\item $N_c=3$ and $\dim\mathrm{U}(1)=1$ from the Radon--Hurwitz numbers $\rho(12)$ and
$\rho(14)$. The same topological invariant that gives 3 spatial dimensions at $d=4$ gives
3 colours at $d=12$.

\item \emph{The gauge window is unique.} $d=12$ is the only dimension in $[5,d_1=19]$ where
$\rho(d)-1=3$ (Theorem~\ref{thm:adams-unique}). The colour count is a forced consequence
of the cascade's geometry, not a selection.

\item $\mathrm{SU}(3)$ and $\mathrm{U}(1)$ are unbroken; $\mathrm{SU}(2)$ is broken.
This follows from two independent arguments: (i) the hairy ball theorem (every tangent
vector field on $S^{12}$ has a zero), and (ii) the Lefschetz theorem
(Theorem~\ref{thm:lefschetz}): no connected compact Lie group can act freely on $S^{12}$.
The Lefschetz version is the stronger group-theoretic statement: SU(2) breaking is an
obstruction on the group action, not merely on a specific field configuration.

\item $m_H/m_W=\pi/2$ from the geodesic distance on $S^{12}$ (0.80\%). The Higgs mass is
the fifth appearance of the cascade's orthogonality axiom.

\item There are exactly three observable charged-fermion generations, from the intersection
of Bott periodicity (which selects Dirac-layer candidates at $d\bmod 8=5$) and the
$d_1=19$ phase transition (which suppresses candidates above $d_1$ by $\sim 289$). The
fourth Bott fermion layer at $d=29$ is a cascade primitive that is suppressed in
charged-fermion channels but sources the heaviest neutrino mass in Part~IVb
(Theorem~\ref{thm:generations}).

\item The gauge window carries the quaternionic phases $\{+1,i,-1\}$, connecting to the
division algebra structure $\mathbb{C}\otimes\mathbb{H}$.

\item The overlap structure of cascade layers is encoded in the lapse function: the Gram
matrix satisfies $G_{d_1 d_2}=N(d_1+d_2)$ (Remark~\ref{rem:gram}).

\item The topological classification of cascade layers---non-trivial-phase layers
($d\bmod 8\in\{5,6\}$) as matter, real-phase Weyl layers ($d\bmod 8=4$) as gauge,
Majorana layers carrying no complex structure---follows from the propagator phase
structure (\cite{paperII}, Corollary~6.8; refined to period~8 via
Clifford in~\cite{paperIII}) and the Bott partition established in this paper.
By Corollary~3.2 of~\cite{paperI}, sphere areas are the only independent cascade
quantities; combining this partition with that corollary provides the topological ingredient
for the matter energy fraction, which follows from the sphere-area fraction of
non-trivial-phase layers in the full cascade tower.
\end{enumerate}

Every result uses the cascade's geometric content (from Papers~I--III) and classical
mathematical theorems (Bott periodicity, Adams, Lefschetz, Poincar\'{e}--Hopf). No physics
beyond the cascade's geometry enters. The quantitative predictions---mass ratios, absolute
masses, coupling constants---are derived in Part~IVb~\cite{paperIVb}.

\begin{thebibliography}{9}
\bibitem{paperI} [Author], \emph{Scale Variance from Orthogonality: How the Unit Ball
Generates $10^{120}$ Orders of Magnitude}, March 2026.
\bibitem{paperII} [Author], \emph{Quantum Mechanics from the Cascade: Effective Theory of
a 4-Dimensional Observer in the Sphere-Area Geometry}, March 2026.
\bibitem{paperIII} [Author], \emph{General Relativity from the Cascade: Lovelock Uniqueness
and the Cosmological Constant of a 4-Dimensional Observer}, March 2026.
\bibitem{paperIVb} [Author], \emph{The Standard Model from the Cascade: Part IVb---Masses,
Couplings, and Precision Predictions from the Geometric-Topological Factorization},
March 2026.
\bibitem{lounesto} P.~Lounesto, \emph{Clifford Algebras and Spinors}, 2nd ed., Cambridge
University Press, 2001.
\bibitem{atiyah} M.~F.~Atiyah, R.~Bott, and A.~Shapiro, ``Clifford modules,''
\emph{Topology} \textbf{3}, 3--38 (1964).
\bibitem{adams} J.~F.~Adams, ``Vector fields on spheres,'' \emph{Annals of Mathematics}
\textbf{75}, 603--632 (1962).
\bibitem{adams60} J.~F.~Adams, ``On the non-existence of elements of Hopf invariant one,''
\emph{Annals of Mathematics} \textbf{72}, 20--104 (1960).
\bibitem{lefschetz} S.~Lefschetz, ``Intersections and transformations of complexes and
manifolds,'' \emph{Transactions of the AMS} \textbf{28}, 1--49 (1926).
\bibitem{poincare} H.~Poincar\'{e}, ``Sur les courbes d\'{e}finies par les \'{e}quations
diff\'{e}rentielles,'' \emph{Journal de Math\'{e}matiques Pures et Appliqu\'{e}es}
\textbf{4}(1), 167--244 (1885).
\bibitem{hopf} H.~Hopf, ``Vektorfelder in $n$-dimensionalen Mannigfaltigkeiten,''
\emph{Mathematische Annalen} \textbf{96}, 225--249 (1927).
\bibitem{kaluza} T.~Kaluza, ``Zum Unit\"{a}tsproblem der Physik,'' \emph{Sitzungsberichte
der Preussischen Akademie der Wissenschaften}, 966--972 (1921).
\bibitem{klein} O.~Klein, ``Quantentheorie und f\"{u}nfdimensionale Relativit\"{a}tstheorie,''
\emph{Zeitschrift f\"{u}r Physik} \textbf{37}, 895--906 (1926).
\bibitem{hurwitz} A.~Hurwitz, ``\"{U}ber die Komposition der quadratischen Formen,''
\emph{Mathematische Annalen} \textbf{88}, 1--25 (1923).
\bibitem{milnor} J.~W.~Milnor, \emph{Topology from the Differentiable Viewpoint},
Princeton University Press, 1965.
\end{thebibliography}

\end{document}
